\section*{A two-background calculation}
\addcontentsline{toc}{section}{A two-background calculation}

The static--dynamic identity can be seen without fitting a large model.  Continue with uniform
three-state references and let the persistent pair component be
\begin{equation}
G_0=
\begin{pmatrix}
2&-1&-1\\
-1&1&0\\
-1&0&1
\end{pmatrix}.
\end{equation}
Introduce a second centered direction
\begin{equation}
H=
\begin{pmatrix}
0&1&-1\\
1&-1&0\\
-1&0&1
\end{pmatrix}.
\end{equation}
Both matrices have zero row and column means, and they are orthogonal under the uniform weighted
inner product.  Let a third position $k$ supply a binary background contrast $C\in\{-1,+1\}$ with
equal probability, and define the effective edge
\begin{equation}
G_{i\leftarrow j}(C)=G_0+\lambda C H.
\label{eq:book-two-background-edge}
\end{equation}

The response of the pair $(i,j)$ now changes when position $k$ changes.  Its mean and residual are
\begin{equation}
\overline G_{i\leftarrow j}=G_0,
\qquad
\delta G_{i\leftarrow j}(C)=\lambda C H.
\end{equation}
Direct calculation gives
\begin{equation}
\|G_0\|_{p_i,p_j}^2=\frac{10}{9},
\qquad
\|H\|_{p_i,p_j}^2=\frac{6}{9}=\frac23.
\end{equation}
Therefore
\begin{equation}
\E_C\|G(C)\|_{p_i,p_j}^2
=\frac{10}{9}+\frac{2\lambda^2}{3},
\end{equation}
and the dynamic fraction is
\begin{equation}
\rho_{i\leftarrow j}^{\mathrm{dyn}}
=\frac{6\lambda^2}{10+6\lambda^2}.
\label{eq:book-two-background-rho}
\end{equation}
At $\lambda=0$ the response is exactly static.  At $\lambda=1$, three eighths of its squared mass
is dynamic.  As $|\lambda|$ grows, the pair is increasingly described by its background modulation
rather than by one persistent block.

The construction also displays interaction order.  The term $CH$ jointly involves the audited
target $i$, source $j$, and background position $k$.  It is stored observationally as a
background-dependent edge, but functionally it is a third-order component.  A scalar graph can
record the variance of this edge across backgrounds; it cannot by itself say which categorical
direction $H$ carries or which position produced the modulation.

Finally, the calculation assumes that $C$ is held fixed while the endpoint categories are varied.
If the coefficient were emitted by a gate that also reads $X_j$, then enumerating the source state
$b$ would change both the message and the gate selecting the message.  Double projection would
still center the resulting table, but Eq.~\ref{eq:book-two-background-edge} would no longer describe
one fixed conditional pair operator.  This is the practical reason endpoint exclusion appears as a
separate requirement rather than as a consequence of gauge fixing.
